\subsection{Background}

\begin{Definition}[{\cite{B.Dubrovin2,B.Dubrovin3}}]
An analytic function $F(t)$, where $t = \bigl(t^1, t^2, \dots, t^n\bigr) \in U \subset \mathbb{C}^n$ defined in an open subset of $\mathbb{C}^n$, is considered a solution of the WDVV (Witten--Dijkgraaf--Verlinde--Verlinde) equation if its third derivatives
\smash{$c_{\alpha\beta\gamma} = \frac{\partial^3 F}{\partial t^{\alpha} \partial t^{\beta} \partial t^{\gamma}}$}
satisfy the following conditions:
\begin{enumerate}\itemsep=0pt
\item[(1)] The coefficients $\eta_{\alpha\beta} = c_{1\alpha\beta}$ form elements of a constant nondegenerate matrix.
\item[(2)] The quantities $c_{\alpha\beta}^{\gamma} = \eta^{\gamma\delta} c_{\alpha\beta\delta}$ represent the structure constants of an associative algebra.
\item[(3)] The function $F(t)$ must be quasi-homogeneous.
\end{enumerate}
\end{Definition}


In \cite[Chapter~1]{B.Dubrovin2}, Dubrovin formulated a geometric interpretation of the WDVV equation which is given by the following.

\begin{Definition}
A Frobenius algebra $\mathcal{A}$ is a finite-dimensional unital, commutative, associative algebra equipped with an invariant, non-degenerate bilinear pairing
$
\eta \colon \mathcal{A} \otimes \mathcal{A} \mapsto \mathbb{C}$,
which is invariant in the following sense
$
\eta(A \bullet B, C) = \eta(A, B \bullet C)$, $ \forall A, B, C \in \mathcal{A}$.
\end{Definition}

\begin{Definition}[{\cite{B.Dubrovin2,B.Dubrovin3}}]
Let $M$ be a complex manifold of dimension $n$. A Dubrovin--Frobenius structure over $M$ consists of the following compatible objects:
\begin{enumerate}\itemsep=0pt
\item[(1)] A family of Frobenius multiplications $\bullet_p \colon T_pM \times T_pM \mapsto T_pM$ that are analytically dependent on $p \in M$. This family induces a Frobenius multiplication on
\begin{gather*}
\bullet \colon\ \Gamma(TM) \times \Gamma(TM) \mapsto \Gamma(TM).
\end{gather*}
\item[(2)] A flat pseudo-Riemannian metric $\eta$ on $\Gamma(TM)$, also known as the Saito metric.
\item[(3)] A unity vector field $e$ that is covariantly constant with respect to the Levi-Civita connection~$\nabla$ for the metric $\eta$, i.e., $\nabla e = 0$.
\item[(4)] Consider the tensor $c(X, Y, Z) := \eta(X \bullet Y, Z)$. We require the 4-tensor
$(\nabla_W c)(X, Y, Z)$
to be symmetric with respect to $X, Y, Z, W \in \Gamma(TM)$.
\item[(5)] An Euler vector field $E$ with the following properties
\begin{gather*}
\nabla \nabla E = 0, \qquad \mathcal{L}_E \eta(X, Y) = (2 - d) \eta(X, Y), \qquad \mathcal{L}_E c(X, Y, Z) = c(X, Y, Z),
\end{gather*}
where $X, Y, Z \in \Gamma(TM)$. Moreover, we require $\nabla E$ to be diagonalizable.
Let $\bigl(t^1, t^2, \dots, t^n\bigr)$ be the flat coordinates with respect to the metric $\eta$. These coordinates are denoted as Saito flat coordinates. The Euler vector $E$ can be explicitly represented as
\begin{gather*}
E = \sum_{i=1}^n ((1 - q_i) t_i + r_i) \partial_i.
\end{gather*}
\end{enumerate}
\end{Definition}

Roughly speaking, Dubrovin--Frobenius manifold is the geometric structure that naturally arise in the domain of any WDVV solution, which is given by a family of Frobenius algebra on the sheaf of holomorphic vector fields, a flat structure and some suitable marked vector fields. An important example of WDVV solutions are the generating function of Gromov--Witten invariants called Gromov--Witten potential. Another source of Dubrovin--Frobenius manifolds comes from Landau--Ginzburg superpotential, which are unfolding of singularities or family of covering over~$\mathbb{CP}^1$. In the analytic theory of Dubrovin--Frobenius manifold, there exist two flat connection. The 1st structure connection is called Dubrovin connection and it is defined below.

\begin{Definition}[{\cite{B.Dubrovin2,B.Dubrovin3}}]
Consider the following deformation of the Levi-Civita connection defined on a Dubrovin--Frobenius manifold $M$,
$
\tilde{\nabla}_{u}v = \nabla_{u}v + zu\bullet v$, $ u,v\in \Gamma(TM)$,
where $\nabla$ represents the Levi-Civita connection of the metric $\eta$, $\bullet$ denotes the Frobenius product, and $z\in \mathbb{CP}^1$. The Dubrovin connection defined in $M\times\mathbb{CP}^1$ is then given by
\begin{gather}
\tilde{\nabla}_{u}v = \nabla_{u}v + zu\bullet v,\qquad
\tilde{\nabla}_{\frac{\rm d}{{\rm d}z}}\frac{\rm d}{{\rm d}z} = 0, \qquad \tilde{\nabla}_{v} \frac{\rm d}{{\rm d}z} = 0,\nonumber\\
\tilde{\nabla}_{\frac{\rm d}{{\rm d}z}}v = \partial_{z}v + E\bullet v - \frac{1}{z}\mu(v).\label{Dubrovin Connection}
\end{gather}
Here, $\mu$ is a diagonal matrix given by
$\mu_{\alpha\beta} = \bigl(q_{\alpha}-\frac{d}{2}\bigr)\delta_{\alpha\beta}$.
\end{Definition}

The deformation of Levi-Civita connection \eqref{Dubrovin Connection} is again a flat connection. In Saito flat coordinates, the Dubrovin connection flat coordinate system, i.e., the solution of
$
\tilde \nabla {\rm d}\tilde t=0$,
 can be written as
 \begin{gather*}
 \bigl(\tilde \nabla_{\alpha} \omega\bigr)_{\beta}=\partial_{\alpha}\omega_{\beta}-zc_{\alpha\beta}^{\gamma}\omega_{\gamma}=0,\\
\bigl(\tilde \nabla_{\frac{\rm d}{{\rm d}z}} \omega\bigr)_{\beta}=\partial_{z}\omega_{\beta}-E^{\sigma}c_{\sigma\beta}^{\gamma}\omega_{\gamma}+\frac{\mu_{\beta}}{z}\omega_{\beta}=0,
\end{gather*}
 where $\omega={\rm d}\tilde t=\omega_{\alpha}{\rm d}t^{\alpha}$ and $\partial_{\alpha}:=\frac{\partial}{\partial t^{\alpha}}$. Alternatively, we could write the Dubrovin flat coordinate system in the matrix form as
\begin{gather}
\partial_{\alpha}\omega=zC_{\alpha}^{\mathsf T}\omega,\qquad
\partial_{z}\omega=\left(\mathcal{U}^{\mathsf T}-\frac{\mu}{z} \right)\omega,\label{Dubrovin connection flat coordinate system in saito coordinates}
\end{gather}
where
\begin{gather}
\omega=\sum_{\beta=1}^n\partial_{\beta}\tilde t {\rm d}t^{\beta}, \qquad C_{\alpha}=\bigl( c_{\alpha\beta}^{\gamma}\bigr), \qquad \mathcal{U}=\bigl( \mathcal{U}_{\beta}^{\gamma} \bigr):=\bigl( E^{\epsilon}c_{\epsilon\beta}^{\gamma} \bigr)\label{gradient flat saito coordinates def}
\end{gather}
or by using the conjugate system
\begin{gather}
\partial_{\alpha}\xi=zC_{\alpha}\xi,\qquad
\partial_{z}\xi=\left(\mathcal{U}+\frac{\mu}{z} \right)\xi,\qquad \text{where}\quad \xi=\eta^{-1}\omega.\label{Dubrovin connection flat coordinate system in saito coordinates conjugate}
\end{gather}

The compatibility of the system \eqref{gradient flat saito coordinates def} is guaranteed by the vanishing of the Riemann tensor associated with the connection \eqref{Dubrovin Connection}.

The connection \eqref{Dubrovin Connection} is more suitable for presenting the Gromov--Witten potential. The second structure connection, also known as the extended Gauss--Manin connection, is more suitable for presenting the Landau--Ginzburg potential. In order to define this connection, we consider the multiplication by the Euler vector field,
\begin{gather}\label{multiplication by the Euler vector field endomorphism}
E\bullet\colon \ \Gamma(TM) \mapsto \Gamma(TM), \qquad X \in \Gamma(TM) \mapsto E\bullet X \in \Gamma(TM).
\end{gather}
Such an endomorphism gives rise to a bilinear form in the sections of the cotangent bundle of~$M$ as follows. Consider $x = \eta(X ), y = \eta(Y ) \in \Gamma(T^*M)$, where $X, Y \in \Gamma(TM)$. An induced Frobenius algebra is defined on $\Gamma(T^*M)$ by
$x \bullet y = \eta(X \bullet Y)$.

\begin{Definition}[{\cite{B.Dubrovin2,B.Dubrovin3}}]
The intersection form is a bilinear pairing in $\Gamma(T^*M)$ defined by ${g^*(\omega_1, \omega_2) = \iota_E(\omega_1 \bullet \omega_2)}$,
where $\omega_1, \omega_2 \in \Gamma(T^*M)$, and $\bullet$ is the induced Frobenius algebra product in $\Gamma(T^*M)$.
\end{Definition}

In the flat coordinates of the Saito metric is given by
\begin{gather}\label{intersection form generic in flat coordinates}
g^{\alpha\beta} = E^{\epsilon}\eta^{\alpha\mu}\eta^{\beta\lambda}c_{\epsilon\mu\lambda}.
\end{gather}

The intersection form $g^{*}$ of a Dubrovin--Frobenius manifold is a flat almost everywhere nondegenerate metric.
The discriminant is the sub manifold such that the intersection form is degenerate
\begin{gather}\label{discriminant}
 \Sigma=\{t\in M\mid \det( g)=0 \} .
\end{gather}

 Hence, the flat coordinate system of the intersection form in Saito flat coordinates
 \begin{gather}\label{Gauss--Manin connection}
 g^{\alpha\epsilon}(t)\frac{\partial^2 x}{\partial t^{\beta}\partial t^{\epsilon}}+\Gamma_{\beta}^{\alpha\epsilon}(t)\frac{\partial x}{\partial t^{\epsilon}}=0
 \end{gather}
 has poles in \eqref{discriminant}, and consequently its solutions $x_a\bigl(t^1,\dots,t^n\bigr)$ are multivalued. The meromorphic connection \eqref{Gauss--Manin connection} is called Gauss--Manin connection of the Dubrovin--Frobenius manifold. The analytical continuation of the solutions $x_a\bigl(t^1,\dots,t^n\bigr)$ has monodromy corresponding to loops around $\Sigma$. This gives rise to a monodromy representation of
 \begin{gather}\label{monodromy on the moduli}
 \pi_1(M\setminus\Sigma)\mapsto {\rm GL}(\mathbb{C}^n),
 \end{gather}
 which is called monodromy of the Dubrovin--Frobenius manifold. Moreover, we can extend the connection \eqref{Gauss--Manin connection} to a connection on $M\times \mathbb{CP}^1$ as follows:
 \begin{gather}\label{Extended Gauss--Manin connection}
\left( g^{\alpha\epsilon}(t)-\lambda\eta^{\alpha\epsilon}\right)\frac{\partial^2 x}{\partial t^{\beta}\partial t^{\epsilon}}+\Gamma_{\beta}^{\alpha\epsilon}(t)\frac{\partial x}{\partial t^{\epsilon}}=0.
 \end{gather}
 The system \eqref{Extended Gauss--Manin connection} is isomonodromic, then its monodromy representation
 \begin{gather}\label{monodromy of GM on the tolal bundle}
 \pi_1(M\times \mathbb{C}\setminus\Sigma_{\lambda})\mapsto {\rm GL}(\mathbb{C}^n), \qquad \text{where} \quad \Sigma_{\lambda}=\det(\lambda-E_{\bullet})
 \end{gather}
is isomorphic to \eqref{monodromy on the moduli}. This fact is straightforward to check once we realise that if $x_{a}\bigl(t^1,t^2,\dots,\allowbreak t^n\bigr)$ is a solution of \eqref{monodromy on the moduli}, then $x_{a}\bigl(t^1-\lambda,t^2,\dots,t^n\bigr)$ is a solution of \eqref{Extended Gauss--Manin connection}. The meromorphic connection \eqref{Extended Gauss--Manin connection} is called extended Gauss--Manin connection.
 Summarising,
 \begin{Definition}[\cite{B.Dubrovin2,B.Dubrovin3}]
 Let $\tilde\nabla$ be the Levi-Civita connection of the intersection form \eqref{multiplication by the Euler vector field endomorphism}. Then
the extended Gauss--Manin connection is connection in $M\times\mathbb{CP}^1$ given by
\begin{gather*}
\tilde\nabla_{\frac{\partial }{\partial t^{\alpha}} }=\frac{\partial }{\partial t^{\alpha}}+(\lambda-E_{\bullet} )^{-1}C_{\alpha}\left(\frac{1}{2}+\mu \right) ,\\
\tilde\nabla_{\frac{\partial }{\partial \lambda} }=\frac{\partial }{\partial \lambda}-(\lambda-E_{\bullet} )^{-1}\left(\frac{1}{2}+\mu \right) .
\end{gather*}
\end{Definition}

The solutions of the flat coordinate system of \eqref{Dubrovin Connection} and \eqref{Extended Gauss--Manin connection} are related by a Fourier--Laplace transform, see Lemma~\ref{Main lemma Mirror symmetry 1}.

\begin{Remark}
The analytic continuation of solutions of the flat coordinates system \eqref{Gauss--Manin connection} and \eqref{Extended Gauss--Manin connection} along any path of $M\setminus{\Sigma}$ and $M\times\mathbb{CP}^1\setminus{\Sigma_{\lambda}}$, respectively,
\begin{gather}
t=\bigl(t^1,t^2,\dots,t^n\bigr)\mapsto (x_1(t),\dots,x_n(t)),\nonumber\\
(\lambda,t)=\bigl(\lambda,t^1,t^2,\dots,t^n\bigr)\mapsto (x_1(\lambda,t),\dots,x_n(\lambda,t))\label{generic period map and extended period map}
\end{gather}
are called period map and extended period map, respectively.
\end{Remark}

 \begin{Remark}
The solutions of flat coordinate system of \eqref{Gauss--Manin connection} and \eqref{Extended Gauss--Manin connection} are quasi-homogeneous in the following sense:
 \begin{gather}\label{quasi homogeneous condition of the solutions xa and xalambda}
 E(x_a(t))=\frac{1-d}{2}x_a(t),\qquad
\left( \lambda\frac{\partial}{\partial \lambda}+E\right)(x_a(\lambda,t))=\frac{1-d}{2}x_a(\lambda,t).
 \end{gather}
 Moreover, let $g^{ab}$ be the coefficients of the intersection form \eqref{intersection form generic in flat coordinates} in its flat coordinates. Then, there is a polynomial relation between Saito flat coordinates $\bigl(t^1,t^2,\dots,t^n\bigr)$ and intersection form flat coordinates $(x_1,x_2,\dots,x_n)$ given by
 \begin{gather}\label{quadratic relation between Saito and Intersection form flat coordinates}
 t_1:=\eta_{1\alpha}t^{\alpha}=g_{ab}x_ax_b, \qquad \text{where}\quad (g_{ab})=\bigl(g^{ab}\bigr)^{-1}.
 \end{gather}
 See in \cite[Appendix~G, Exercise~G.1]{B.Dubrovin2}.
 \end{Remark}

Recall that a point in a Dubrovin--Frobenius manifold is called semisimple if the Frobenius algebra in $T_pM$ is semisimple. It is worth noting that semisimplicity constitutes an open condition. The Dubrovin--Frobenius structure around semisimple points becomes rather simple. Specifically, the Frobenius algebra becomes trivial. Moreover, both the Saito metric and the endomorphism resulting from the multiplication by the Euler vector field~\eqref{multiplication by the Euler vector field endomorphism} are diagonal around such a point. %More precisely,

\begin{Proposition}[\cite{B.Dubrovin2,B.Dubrovin3}]
Let $(u_1,u_2,\dots,u_n)$ be pairwise distinct roots of the characteristic equation
\begin{gather}\label{spectral curve}
\det\bigl(g^{\alpha\beta} -u\eta^{\alpha\beta} \bigr) = 0.
\end{gather}
Then, the relation $(u_1(t),u_2(t),\dots,u_n(t))$ can serve as local coordinates, which are called canonical coordinates. In these coordinates, the Frobenius multiplication, Saito metric, unit vector field and Euler vector field can be written as
\begin{gather}
\frac{\partial}{\partial u_i}\bullet \frac{\partial}{\partial u_j}=\delta_{ij}\frac{\partial}{\partial u_i},\qquad \eta=\sum_{i=1}^{n}\psi_{i1}^2 ({\rm d}u_i)^2, \qquad e=\sum_{i=1}^n\frac{\partial}{\partial u_i},\qquad E=\sum_{i=1}^n u_i\frac{\partial}{\partial u_i},\label{Dubrovin Frobenius structure in canonical coordinates 1}
\end{gather}
where the matrix $\Psi=( \psi_{i\alpha})$ is given by
$
\psi_{i\alpha}=\psi_{i1}\frac{\partial u_i}{\partial t^{\alpha}}$.
\end{Proposition}
At this stage, we can define a Landau--Ginzburg superpotental, which can be found in \cite[Definition~5.7]{B.Dubrovin3}.

\begin{Definition}[\cite{B.Dubrovin3}]\label{Landau--Ginzburg superpotential definition}
Let $D$ be an open domain of a Riemann surface. A Landau--Ginzburg superpotential associated with a Dubrovin--Frobenius manifold $M$ of dimension $n$ consists of a~function $\lambda(p,u)$ on $D \times M$ and an Abelian differential $\phi$ in $D$ satisfying
\begin{itemize}\itemsep=0pt
\item The critical values of $\lambda(p,u)$ are the canonical coordinates $(u_1,\dots,u_n)$. In other words, the canonical coordinates $(u_1,u_2,\dots ,u_n)$ are defined by the following system
$\lambda(p_i)=u_i$, $\frac{{\rm d} \lambda}{{\rm d}p}( p_i)=0$.
\item For some cycles $Z_1,\dots,Z_n$ in $D$ the integrals
\begin{gather}\label{Mirror symmetry}
\tilde t_j(z,u)=\frac{1}{z^{\frac{3}{2}}}\int_{Z_j}{\rm e}^{z\lambda(p)}\phi, \qquad j=1,\dots,n,
\end{gather}
converges and give a system of independent flat coordinates for the Dubrovin connection~$\tilde\nabla$ in canonical coordinates, i.e., the matrix
$Y=\Psi\eta^{-1}\omega:=\bigl(\psi_{i\alpha}\eta^{\alpha\beta}\partial_{\beta}\tilde t \bigr)$
is a solution of the following system
\begin{gather}\label{Dubrovin connection in canonical coordinates 1}
\frac{\partial Y}{\partial u_i}=( zE_i+V_i )Y,\qquad
\frac{{\rm d}Y}{\partial z}=\left( U+\frac{V}{z} \right)Y,
\end{gather}
where
\begin{gather}\label{U,V, Vi in canonical coordinates}
U=\Psi\mathcal{U}\Psi^{-1}, \qquad V=\Psi\mu\Psi^{-1}, \qquad E_i=( \delta_{ij}\delta_{ik} ), \qquad V_i:=\frac{\partial \Psi}{\partial u_i}\Psi^{-1}.
\end{gather}

\item
The following expressions for the coefficients of the tensors Saito metric $\eta$, intersection form $g^{*}$ and the structure constants $c$, in any coordinate system $(x_1,x_2,\dots,x_n)$, holds true
\begin{gather}
\eta_{ij}=\sum \res_{{\rm d}\lambda=0} \frac{ \partial_i\lambda \partial_j\lambda }{d_p\lambda }\phi, \qquad
g_{ij}=\sum \res_{{\rm d}\lambda=0} \frac{ \partial_i\log\lambda \partial_j\log\lambda }{d_p\log\lambda } \phi,\nonumber\\
c_{ijk}=\sum \res_{{\rm d}\lambda=0} \frac{ \partial_i\lambda \partial_j\lambda \partial_k\lambda }{d_p\lambda }\phi.\label{residue expression for eta, intersection form and structure constants}
\end{gather}
\end{itemize}
\end{Definition}

\subsection{Problem setting}

Let $N_k$ be the number of rational curves $\mathbb{CP}^1 \rightarrow \mathbb{CP}^2$ of degree $k$ passing through $3k-1$ generic points. Kontsevich \cite{M.Kontsevich} proved that the generating function
\begin{gather}\label{generating function of Gromov Witten CP2}
F\bigl(t^1, t^2, t^3\bigr) = \frac{\bigl(t^1\bigr)^2t^3}{2} + \frac{\bigl(t^2\bigr)^2t^1}{2} + \sum_{k=0}^{\infty} \frac{N_k}{(3k-1)!}{\rm e}^{kt^2}\bigl(t^3\bigr)^{3k-1}
\end{gather}
satisfies a WDVV equation. Moreover, in \cite{M.Kontsevich} the Gromov--Witten invariants $N_k$ were computed as follows. Note that the associativity equation for any 3d Dubrovin--Frobenius manifold is given~by
\begin{gather}\label{associativity equation generic}
c_{223}^2 = c_{333} + c_{222}c_{233}.
\end{gather}
Defining the function
\begin{gather}\label{main Gromov Witten potential}
\Phi(X) = \sum_{k=0}^{\infty} \frac{N_k}{(3k-1)!}{\rm e}^{kX}, \qquad X := t^2 + 3\ln t^3,
\end{gather}
we obtain the following third-order differential equation for the function $\Phi(X)$ by substituting~\eqref{generating function of Gromov Witten CP2} and \eqref{main Gromov Witten potential} into \eqref{associativity equation generic},
\begin{gather}\label{associativity equation}
-6\Phi + 33\Phi' - 54\Phi'' - \bigl(\Phi''\bigr)^2 + \Phi''' \bigl( 27 + 2\Phi' - 3\Phi'' \bigr) = 0.
\end{gather}
Hence, using the Taylor series expansion of \eqref{main Gromov Witten potential} in \eqref{associativity equation}, we obtain the following recursion:
\begin{gather*}
N_d = \sum_{m=1}^{d-1} \left[ {{3d-4}\choose{3m-2}} m^2(d-m)^2 - {{3d-4}\choose{3m-3}} m(d-m)^3 \right]N_mN_{d-m},
\end{gather*}
which explicitly provides all Gromov--Witten $N_k$ by setting $N_1=1$. In principle, the Gromov--Witten potential \eqref{main Gromov Witten potential} is only a formal power series. However, in \cite{DiFrancesco} Di Francesco and Itzykson derived the following asymptotic behavior for~$N_k$:
\begin{gather*}
\frac{N_k}{(3k-1)!} = ba^kk^{-\frac{7}{2}} \left(1 + O\left(\frac{1}{k}\right)\right),
\end{gather*}
where $a$ and $b$ are numerically estimated as
$a = 0.138$, $ b = 6.1$,
which implies that the radius of convergence of \eqref{main Gromov Witten potential} is given by $\frac{1}{a}$. Then, the analytic domain of \eqref{generating function of Gromov Witten CP2} is given by
\begin{gather}\label{WDVV domain}
D = \left\{\bigl(t^1, t^2, t^3\bigr) \in \mathbb{C}^3 \mid \big|{\rm e}^{t^2}\bigl(t^3\bigr)^3\big| < \frac{1}{a}\right\}.
\end{gather}
Furthermore, the Euler vector field has the following form
$E = t^1\partial_1 + 3\partial_2 - t^3\partial_3$.

Our aim is to investigate the analytic properties of \eqref{generating function of Gromov Witten CP2} by constructing the corresponding 1D Landau--Ginzburg superpotential from the Gauss--Manin connection of $QH^{*}\bigl(\mathbb{CP}^2\bigr)$.

